\paragraph{A completed improvement from the infinite prime-density tail.}
The first 1000 periods need not be the endpoint of the density calculation.
For $0\le a<b\le1$, the positive function
$f(t)=(t+a)^{-1}-(t+b)^{-1}$ is decreasing and convex. The trapezoid and
midpoint inequalities give the explicit two-sided tail bounds
\[
 \log\frac{K+b}{K+a}+\frac12\left(\frac1{K+a}-\frac1{K+b}\right)
 \le \sum_{t=K}^{\infty}f(t)
 \le \log\frac{K-\frac12+b}{K-\frac12+a}.
\]
The logarithms are enclosed by the positive series
$\log z=2\sum_{j\ge0}u^{2j+1}/(2j+1)$, where $u=(z-1)/(z+1)$,
with its geometric remainder bound. Applying this with $K=1000$ to
every baseline cell and every one of the fourteen established periodic
shells yields
\[
 29.9127079828158127 <\gamma_{\rm established}
 <29.9127079831445955.
\]
An independent Arb evaluation using digamma functions gives
$\gamma_{\rm established}=29.912707983034852259666192\ldots$.
Using only the safe rational value $\gamma=29.91270798$, and the same
accepted analytic rates as before, gives
\[
\begin{array}{c|cc|c}
 &\sigma&\tau&\liminf\delta_N\\\hline
 B=0&92.538828585866&20.268710661589&>0.21902925\\
 B\ne0&111.538828585866&1.268710661589&>0.01137460.
\end{array}
\]
This small improvement uses the already established divisibility statement;
it does not rely on the two additional numerical windows.

\paragraph{The larger improvement identified by the residual gcds.}
The new second-order lemma establishes an additional common prime factor on
\[
 \{n/p\}\in(1/10,2/19)\ \cup\ (11/12,15/16).
\]
They pass all 843 eligible prime instances at
$n=49,99,249,499,999,1999,3999,7999$, for both readouts.
Their complete prime densities lie respectively in
\[
 (0.50751800694358,0.50751800694492),\qquad
 (0.03832125404609,0.03832125405130).
\]
Using their established uniform divisibility, the same rate
calculation gives $\gamma>30.45854724380549$, and consequently
$\liminf\delta_N>0.22626234$ for $B=0$ and
$\liminf\delta_N>0.01634832$ for $B\ne0$.
These bounds use the separately established uniform second-order trace
lemma and the accepted analytic estimates. The finite checks alone do not
establish the arithmetic lemma.
